\documentclass[12pt]{article}
\usepackage{graphicx} % Required for inserting images
\usepackage{setspace}
\usepackage{fancyhdr}
\usepackage{amsmath}
\usepackage{caption}
\usepackage{graphicx}
\usepackage{booktabs}
\usepackage{adjustbox}
\usepackage{url}
\usepackage[utf8]{inputenc}
\usepackage{csquotes}
\usepackage{natbib}
\usepackage{pdfpages}
\usepackage[raggedright]{ragged2e}
\usepackage[hidelinks]{hyperref}
\usepackage[letterpaper,margin=1in]{geometry}
\bibliographystyle{apsr}   % or plainnat, chicago, apa, etc.
\setcitestyle{authoryear,open={(},close={)},aysep={ }}

\title{We Don't Do That: National Identity and Nuclear Non-Use} 
\author{Wyatt King}
\date{Draft: February 27}

\begin{document}

\maketitle

\section{Introduction}

In international relations scholarship, political scientists have often taken after microeconomic theory in assuming that states are rational and selfish. Neorealist scholars presuppose that states are strictly security seeking, unmoved by considerations of fairness or altruism in an international system predicated on anarchy and self-help. Since the “constructivist turn,” however, scholars have challenged the descriptive accuracy of assuming states to be egoistic. Instead, individual leaders, domestic constituencies, and states as units in and of themselves can be constrained by norms of how they are meant to behave. Crucial to how norms are theorized to affect state behavior is identity. If an action is so morally reprehensible as to contradict a state's perception of itself, then it will refrain from whatever that action is.

d

In political behavior research outside of international relations, scholars have found priming national identity to be salient in reducing affective polarization, changing opinions toward immigration, and \_\_. Equivalent experimentally-driven research have yet to be applied to scenarios urgent to  \\

Rather than estimating average causal effects in a similar manner 

As conventionally understood by constructivists, norms are an “expectation of appropriate behavior for states with a given social identity.” A state will only follow a given course of action if it coheres with its sense of self. If an action is so morally egregious as to contradict how a state understands itself and its interests, then it will refrain from doing said action. Among the most robust norms in the international system, according to constructivist scholars, is that against nuclear weapons use. Although the “nuclear taboo” as Nina Tannenwald has called it is both the product of

Among the mechanisms through which norms are theorized to affect policy preferences is national identity. In international relations, constructivist scholars contend that states are sensitive to norms because they seek to validate social identities. States are not strictly security-seekers, as neorealists suggest, but in fact, have a specific conception of how they are meant to behave and what principles they are supposed to uphold. As a result, states' behavior may be limited by their self-perceived conception of who they are. 

International relations scholars have tended to make this assumption most strongly, suggesting that states are strictly security seeking


However, scholars across subfields are increasingly criticizing this assumption. Rather than being motivated by selfishness, they argue that actors can be other-regarding in given circumstances. Among the mechanisms

International relations scholars tend to be among the most prominent scholars making this assumption, casting states as being in a zero-sum competition for power and security. Across disciplines within political science, however, scholars have increasingly criticized this notion, suggesting that actors may possess altruistic motivations in given contexts.

In his Nobel Prize lecture, Thomas \citet{schelling_astonishing_2006} identified “the most spectacular event” of the late twentieth century as one that did not occur. “We have enjoyed sixty years without nuclear weapons exploded in anger,” he said, “Can we make it through another half dozen decades?”

At the time Schelling gave his speech, burgeoning scholarship argued that a taboo had developed around nuclear weapons during the Cold War, making their use “unthinkable” by its end \citep{tannenwald_stigmatizing_2005}. While scholars of the taboo primarily couched their argument in elite-level case studies, they extended it to public opinion, suggesting that it too exhibited a widespread aversion to the use of nuclear weapons. Over the last ten years, a new generation of scholarship on the taboo has challenged this view by fielding survey experiments in which respondents are asked how strongly they would support the first-use of nuclear weapons in different hypothetical scenarios. Consistently, these surveys have found that a majority of respondents in Western, nuclear-armed states would support a first-strike nuclear attack \citep{sagan_revisiting_2017, haworth_what_2019,allison_under_2022, dill_kettles_2022}.

Although these surveys provide robust evidence that the public's aversion to the use of nuclear weapons is not absolute, they do not dismiss that there is a norm against it. All else being equal, people prefer using conventional weapons to nuclear ones and sizable portions of public opinion still strongly oppose using nuclear weapons even when they would be strategically useful \citep{press_atomic_2013}. Consequently, if it is true that the public possesses a normative aversion to the use of nuclear weapons, then what factors can reify this norm and expand its effect?

To answer this question, I return to initial theorizing on the nuclear taboo and argue that people's conceptions of their state's national identity may be one such factor. Following from constructivist theory, I assume that states pursue patterns of behavior that cohere with their self-perceived identities \citep{tannenwald_nuclear_2007,wendt_social_1999}. Crucially, these identities can be constructed such that they disqualify certain actions as being beyond the pale--so opposed to a state's national identity as to be antithetical to its interests.

To test whether norms have such “constitutive effects,” I field an online survey experiment of Indian citizens similar to those pioneered by Scott Sagan and Benjamin Valentino. Exposing respondents to a hypothetical scenario in which India launches a preemptive counterforce attack against Pakistan, I test how national identity-based arguments opposing nuclear use change the degree of public approval for the attack. Furthermore, considering that the institutionalization of norms may affirm people's conceptions of national identity as being opposed to nuclear use, I gauge how informing respondents of India's No First Use policy changes public support for nuclear use both independent of national identity based arguments and in conjunction with them.

Ultimately, the purpose of this thesis is twofold. Narrowly, it contributes to the upswell of survey experiments from the past ten years, in which scholars have sought to understand what factors impact public support for the use of nuclear weapons. It does so by interrogating the effects that national identity and institutionalization can have, and whether those effects are significant in the manner that scholars of the taboo suggest. More broadly, it seeks to evaluate claims in constructivist literature that norms can have a “constitutive effect,“ in which people's foreign policy preferences are impacted by their impression of state identity.

\section{Case Selection}
To test whether national identity can alter support for nuclear first use, I field a survey experiment modeled off of those created by Sagan and Valentino. However, instead of fielding it in a Western nuclear armed state, as most survey experiments up to now have done, I conduct it in India. I do so because India has a lasting tradition of political elites arguing that non-violence is politically pragmatic, morally obligatory, and part of Indian national identity. As a result, India amounts to a most likely case in which national identity-based arguments could mitigate support for nuclear weapons use.

The incorporation of non-violence into Indian national identity stems from the nationalist movement, in which political elites argued that non-violence was both an effective anti-imperial strategy and a moral obligation. None did so as clearly as Mahatma Gandhi. According to Gandhi, people were spiritually required to uphold \textit{ahimsa} (a religiously-inflected understanding of non-violence shared by Hindus, Jains, and Buddhists that literally means “to not kill or injure”). Gandhi's understanding of \textit{ahimsa} was deontological and made no meaningful distinction between offensive and defensive wars. Following from his categorical aversion to violence, Gandhi's opposition to nuclear use was staunch and unequivocal. Reflecting on the American bombings of Hiroshima and Nagasaki, he said that, “Unless now the world adopts non-violence, it will spell certain suicide for mankind.” When challenged that nuclear weapons could be an instrument to achieving peace by creating stable deterrence, Gandhi was incisive. He compared nuclear weapons' deterrent effect to a “a man glutting himself with dainties.” States will refrain from violence “only to return with redoubled zeal after the effect of nausea.”

Following Gandhi's death, his peers in the Indian nationalist movement adapted non-violence to Indian foreign policy, at least rhetorically. In 1954, India's Prime Minister Jawaharlal Nehru became the first world leader to call for disarmament. During negotiations over what would become the Nuclear Non-Proliferation Treaty, India sought to give teeth to these demands. Criticizing nuclear-armed powers for their hypocrisy in wanting to stop proliferation while allowing themselves to maintain and enlarge their nuclear arsenals, India refused to sign a multilateral non-proliferation treaty unless nuclear weapons states were required to disarm. Otherwise, the NPT amounted to what Indian policymakers called “nuclear apartheid,” in which wealthy states would secure their access to nuclear weapons and protect a crucial strategic advantage over poorer, primarily post-colonial, states. When the U.S., Soviet Union, and others rejected calls to disarm, India refused to join the NPT.

To be sure, for however much India's opposition to the NPT stemmed from sincere moral feeling, it was also cut with an understanding of strategic necessity. In 1964, China conducted its first nuclear tests and proceeded to build up its arsenal. Due to long-standing border disputes between the two countries, which culminated in the 1962 Sino-Indian War, China's possession of nuclear weapons risked tilting the balance of power in the Himalayas in its favor. By not signing the NPT, India guaranteed itself the option to later balance against China by testing nuclear weapons—as it would do in 1974. 

Even as India pursued nuclear weapons, it continued to frame itself as a non-violent and responsible power. The codename for India's nuclear test was “smiling Buddha,” and in announcing its success, Indira Gandhi deemed it a “peaceful nuclear explosion.” Aside from naming conventions that could amount to cheap talk, India did not press its nuclear advantage over Pakistan after testing. It did not invest in ballistic missile systems or submarines that would make its deterrent credible, and it did not test nuclear weapons again until more than two decades later. Moreover, India continued its advocacy for global disarmament. Indira Gandhi coordinated a group of five other states to advocate for disarmament, and in 1988, her son, Prime Minister Rajiv Gandhi, proposed a plan at a special U.N. summit on disarmament that would eliminate nuclear weapons globally by 2010. 

Compared to Congress-led governments of the mid- to late-twentieth century, the Bharatiya Janata Party has been far less averse to nuclear weapons. With the BJP leading the government for the first time, Prime Minister Atal Vajpayee authorized India's second nuclear weapons tests in 1998, despite security analysts widely believing that doing so unnecessarily inflamed tensions in South Asia. After Pakistan tested nuclear weapons in response, the two countries briefly went to war in 1999. In the wake of the Kargil War, Vajpayee announced that India would follow a No First Use policy, in which it would refrain from using nuclear weapons unless first attacked by biological, chemical, or nuclear weapons. 

Today, there is widespread doubt regarding whether India would actually adhere to its NFU during a crisis with Pakistan \citep{clary_indias_2019}. In its 2014 election manifesto, the BJP said that India's NFU should be “revised and updated.” When asked by Reuters journalists what this revision entailed, anonymous BJP sources said cryptically to “read the manifesto.” After public backlash, Narendra Modi clarified that India would maintain its NFU, but it has continued to take action that raise doubts regarding the legitimacy of these claims. In 2019, after a terrorist attack in India-controlled Kashmir killed \_\_ civilians, India struck non-disputed Pakistani territory for the first time since 1971. Then-U.S. Secretary of State Mike Pompeo claimed in his memoir that had the U.S. not been involved, the conflagration risked spiraling and leading to nuclear use. Several months after the attack Indian Defense Minister Rajnath Singh, standing at the site of India's second nuclear tests, announced that although India had previously “strictly adhered” to its NFU, “what happens in future depends on the circumstances.” In May 2025, India and Pakistan engaged in a skirmish similar to that in 2019, during which India struck a Pakistani air base near the military headquarters responsible for overseeing its nuclear arsenal. The crisis was also the first time that India used cruise missiles against Pakistan, and Pakistan used ballistic missiles against India. India has also built out its intelligence gathering operations, improved the range of its ballistic missile system, and developed multiple targetable re-entry vehicles that would enable it to launch a counterforce attack against Pakistan that would disable its nuclear weapons facilities.

Even as public officials have raised doubt about the legitimacy of India's NFU, they have clung to the language of India being a “responsible” nuclear state. In the same speech where he questioned that India would abide by its NFU, Singh said that “India attaining the status of a responsible nuclear nation is a matter of national pride.” After public outcry over the BJP's election manifesto, Modi called India's NFU “a reflection of cultural inheritance.” The contradictory nature of these statements suggests that although Indian elites may be disillusioned with the utility of India's NFU, they still believe that it is a useful signaling device either to domestic political suspicious of an aggressive nuclear posture or adversaries skeptical of Indian nuclear use.

\bibliography{references}

\end{document}